\section{Capítulo~\ref{sec:sensors}. Las entradas de la máquina}

La estructura de los sensores se ha medido directamente: con registros de corriente de células individuales, de las vibraciones de la cóclea y con recuentos de células de la retina; el límite de sensibilidad y la fórmula del ruido de disparo se deducen de la física, y que los códigos de los sensores estén ajustados para la máxima información es una hipótesis con un apoyo sólido.

{\footnotesize
\setlength{\tabcolsep}{3pt}\renewcommand{\arraystretch}{1.15}
\begin{longtable}{>{\raggedright}p{0.5cm}>{\raggedright}p{4.0cm}>{\raggedright}p{5.6cm}>{\raggedright}p{2.3cm}>{\raggedright\arraybackslash}p{1.7cm}}
\toprule
N.º & Afirmación & Experimento y qué se midió & Fuente & Estado \\
\midrule\endfirsthead
\toprule
N.º & Afirmación & Experimento y qué se midió & Fuente & Estado \\
\midrule\endhead
1 & El sensor es un transductor: el receptor produce una corriente, y la salida de las células sensoriales del ojo y del oído es glutamato hacia el bus \nt{GLUT}; las primeras células no emiten espigas (fig.~\ref{fig:transducer}) & Los bastones y conos de la tortuga liberan un aminoácido excitador: lo detectan canales de glutamato en un parche de membrana escindido. Células ciliadas internas de la cóclea de rata: las corrientes en las terminales de las fibras nerviosas pasan por receptores AMPA de glutamato, y la frecuencia de liberación crece con la despolarización gradual de la célula hasta $150$\,Hz & \cite{CopenhagenJahr1989,GlowatzkiFuchs2002} & medido \\
2 & El sensor de luz en la oscuridad responde a un solo fotón con una corriente de alrededor de un picoamperio & Electrodo de succión sobre el segmento externo de un bastón de sapo: las respuestas a destellos tenues están cuantizadas, evento de $\approx1$\,pA ($5\%$ de la corriente oscura) con dispersión de $0.2$\,pA, eficiencia de $0.5$. Las personas informan de la llegada de un solo fotón más a menudo que al azar & \cite{Baylor1979,Tinsley2016} & medido \\
3 & El sensor de sonido detecta desplazamientos de menos de un nanómetro & Método Mössbauer, membrana basilar de la cóclea del cobayo en el punto de $16$--$19$\,kHz: en el umbral de respuesta del nervio auditivo, la velocidad de la membrana es de unos $0.04$\,mm/s, es decir, un desplazamiento de $\approx0.35$\,nm (cálculo nuestro). Se midió la membrana, no el haz del sensor & \cite{Sellick1982,Hudspeth1989} & medido \\
4 & En la oscuridad la precisión está limitada por el ruido de disparo de los fotones~\eqref{eq:photonnoise}; con luz débil la acumulación es más larga & La fórmula se deduce de la estadística de Poisson. En el bastón, el ruido con luz tenue coincide con la suma de eventos cuánticos independientes; sobre un fondo brillante, la respuesta a un destello es menor y más breve. La cifra de “décimas de segundo” no se comprobó aquí por separado & \cite{Baylor1979} & teoría \\
5 & El rango de entrada es de diez órdenes para la luz y doce para el sonido, y el de la frecuencia de espigas, de menos de tres órdenes & Para el sonido: la respuesta de la membrana basilar a la frecuencia característica crece con una pendiente de hasta $0.2$\,dB/dB en la banda de $40$--$80$\,dB, y la compresión se ve también por encima de $100$\,dB. Los números de órdenes son estimaciones estándar, sin fuente en el capítulo & \cite{Ruggero1997} & medido \\
6 & Respuesta del sensor~\eqref{eq:nakarushton}, con $\sigma$ siguiendo el brillo medio y desplazando la curva por el eje logarítmico (fig.~\ref{fig:adaptation}) & La fórmula se ajustó por primera vez a los potenciales S de la retina de pez. Conos de tortuga: las respuestas obedecen $V=I V_m/(I+\sigma)$ y abarcan $\sim3.5$ órdenes de brillo; tras $2$--$3$ minutos de fondo, la curva se desplaza horizontalmente sin cambiar de ancho. La relación con la división por la actividad total es una revisión & \cite{NakaRushton1966,NormannPerlman1979,Carandini2012} & medido \\
7 & Los sensores transmiten cambios, y sus códigos están ajustados para la máxima información por espiga (proposición~\ref{prop:sensors}) & El principio se propuso teóricamente. La prueba más fuerte: en la mosca, la característica “contraste $\to$ respuesta” de las neuronas de la primera capa coincide con la distribución acumulada de contrastes de las escenas naturales, y así se alcanza la máxima información & \cite{Barlow1961,Laughlin1981} & hipótesis \\
8 & Los sensores de encendido y apagado son un código de dos cables; la cámara de eventos lo repite en silicio & Revisión de experimentos: los canales de encendido y de apagado de la retina se separan ya en la primera sinapsis y pueden anularse por separado. Cámara: $128\times128$ píxeles sin fotogramas, rango de $120$\,dB, latencia de $15$\,\textmu s & \cite{Schiller1992,Lichtsteiner2008,Gallego2022} & medido \\
9 & En el ojo hay $\sim10^8$ sensores de luz y $\sim10^6$ neuronas de salida: compresión de $\sim100$ veces & Recuento en retinas humanas completas: en promedio $92$ millones de bastones y $4.6$ millones de conos; células de salida (ganglionares), de $0.7$ a $1.5$ millones según el ojo & \cite{Curcio1990,CurcioAllen1990} & medido \\
10 & En el ratón, el ojo tiene más de treinta canales de salida & Ratón: imagen de calcio con dos fotones de más de $11\,000$ células de salida y agrupamiento de sus respuestas: bastante más de $30$ tipos funcionales. En humanos el número no se ha medido & \cite{Baden2016} & medido \\
11 & El ojo del cobayo transmite $\sim875\,000$ bit/s; el cálculo para el ojo humano da $\sim10^7$ bit/s & Cobayo, matriz multielectrodo, movimiento en escenas naturales: $6$--$13$ bit/s por célula, $\sim875\,000$ bit/s para $\sim10^5$ células de salida. Para el ser humano ($\sim10^6$ células), $\sim10^7$ bit/s es un cálculo de los autores & \cite{Koch2006} & medido \\
12 & El oído es un banco de filtros pasa banda con un mapa de frecuencias casi logarítmico (fig.~\ref{fig:filterbank}) & El lugar de máxima vibración de la membrana basilar depende de la frecuencia; la función “frecuencia--lugar” es casi exponencial y describe con una sola forma las cócleas humanas y de animales vivos & \cite{Greenwood1990,Robles2001} & medido \\
13 & Los resonadores son activos: parte de los sensores amplifica las vibraciones débiles miles de veces en amplitud ($66$--$76$\,dB), y la ganancia cae al aumentar el volumen & Vibrometría láser de la base de la cóclea de la chinchilla: con tonos débiles a la frecuencia característica, ganancia de $66$--$76$\,dB respecto del estribo; la muerte reduce la sensibilidad en $60$--$81$\,dB ($10^3$--$10^4$ veces en amplitud) y elimina la compresión. La fuente de fuerza son las células ciliadas externas & \cite{Ruggero1997,Robles2001} & medido \\
14 & El amplificador del oído está al borde de la autooscilación & Cálculo: cerca de un punto de bifurcación de Hopf son inevitables la compresión del rango, la sintonía aguda y los tonos de combinación, todas las no linealidades conocidas de la audición. Que la cóclea esté ajustada así no se ha demostrado directamente & \cite{Eguiluz2000} & hipótesis \\
15 & A bajas frecuencias las espigas van en fase con el sonido (en el cobayo el enganche se debilita desde $\sim600$\,Hz y se nota hasta $\sim3.5$\,kHz), y con ellas se halla la dirección; a altas queda el código de lugar & Nervio auditivo del cobayo: el enganche de fase empieza a caer hacia $600$\,Hz y desaparece por encima de $3.5$\,kHz; en el gato y el mono el límite está aproximadamente una octava más arriba. La diferencia de tiempo de llegada a los dos oídos: revisión & \cite{PalmerRussell1986,Grothe2010} & medido \\
16 & Los demás sensores (tabla~\ref{tab:sensors}): el olfato tiene $\sim400$ tipos de receptores, uno por neurona, y un código combinatorio; el gusto funciona en parte con ATP y serotonina; el tacto, con sensores de adaptación rápida y lenta; calor y frío, por separado & Ratón: un receptor reconoce muchos olores, un olor activa muchos receptores, y olores distintos activan conjuntos distintos. Sin receptores de ATP el nervio del gusto no responde; los botones gustativos liberan serotonina. Registro de fibras individuales de la piel de la mano humana: cuatro tipos según la velocidad de adaptación. El canal del frío es distinto de los del calor & \cite{Mombaerts2004,Malnic1999,Finger2005,Huang2005,JohanssonVallbo1979,McKemy2002} & medido \\
\bottomrule
\end{longtable}}
